\scanpage{71}
Supposons $X$ contenu dans un schéma lisse $Y$.
\deleted{la conjecture ci-dessus est vraie pour les $Y(i)^{\mathrm{an}}$ ($Y(i)=$ $i$ième voisinage infinitésimal de $X$ sur $Y$)} alors l'hom. \deleted{\unsure{on a}}\inserted{canonique}
\[
H^*(X_{\mathrm{cris}},\mathcal O_{X_{\mathrm{cris}}})\longrightarrow H^*(X_{\mathrm{strat}},\mathcal O_{X_{\mathrm{strat}}})
\]
est également un isom.

Le premier énoncé résulte de l'isom de comparaison plus haut, et de
\[
H^*(X^{\mathrm{an}},\mathcal F_{X^{\mathrm{an}}}^{\deleted{\unsure{\mathrm{an}}}*})\simeq H^*(X^{\mathrm{an}},\mathbf C).
\]
Pour le deuxième, les développements du n° 2 donnent aussitôt
\[
H^*_{\mathrm{cris}}(X,\mathcal O_{X_{\mathrm{cris}}})=\varprojlim_i H^*_{\mathrm{cris}}(Y(i),\mathcal O_{Y(i)_{\mathrm{cris}}})
\]
Or d'après la première partie, les hom de transition sont des isom, et on gagne.

Cela donne donc des espoirs à la conjecture de l'exposé précédent.

\uline{Pb} Supposons de façon générale $X$ \uuline{plat}\inserted{(loc. de prés. finie)} sur $S$, et $F$\inserted{sur $X_{\mathrm{cris}}$}, satisfaisant aux hyp. habituelles (p.ex. provenant d'un Module qu. coh. stratifié) À prouver qu'alors
\[
H^*(X_{\mathrm{cris}},\deleted{\mathcal O_{X_{\mathrm{cris}}}}F)\xrightarrow{\sim}H^*(X_{\mathrm{strat}},F\ (\text{site stratifié})).
\]
\mspage{70}
